%
\subsection*{Un paso de Euler, Heun y RK4}
\label{CH03-B01}
Para $y'=f(t,y)$, $t_{n+1}=t_n+h$, Euler usa
$y_{n+1}=y_n+h f(t_n,y_n)$. Heun calcula un predictor y promedia dos
pendientes:
\[
 k_1=f(t_n,y_n),\quad \widetilde y=y_n+hk_1,\quad
 k_2=f(t_n+h,\widetilde y),\quad
 y_{n+1}=y_n+\tfrac h2(k_1+k_2).
\]
En RK4, con la convención de que $k_i$ es una pendiente,
\[
\begin{aligned}
 k_1&=f(t_n,y_n),& k_2&=f(t_n+h/2,y_n+hk_1/2),\\
 k_3&=f(t_n+h/2,y_n+hk_2/2),&
 k_4&=f(t_n+h,y_n+hk_3),\\
 y_{n+1}&=y_n+\tfrac h6(k_1+2k_2+2k_3+k_4).
\end{aligned}
\]
Cada evaluación tiene un tiempo y un estado propios. La pendiente no es
el incremento $hk_i$. Compárense resultados en el mismo instante y,
si se compara coste, cuéntense las evaluaciones de $f$; redondéese al final.
Un caso particular no establece una superioridad universal.
Consultas: soluciones de \exref{M2-E0323} y \exref{M2-E0298}.

%
\subsection*{Error, estabilidad y restricciones del paso}
\label{CH03-B02}
Frente a un valor de referencia $y_*$, el error absoluto es
$|y_n-y_*|$; el relativo es $|y_n-y_*|/|y_*|$ cuando $y_*\ne0$.
En $y'=-ay$, $a>0$, Euler da $y_{n+1}=(1-ah)y_n$.
El módulo del multiplicador controla la amplificación; su signo afecta
a la no negatividad. Ninguna de esas propiedades garantiza precisión.
Para conservar un intervalo $I$, compruébese que el mapa de un paso lleva
\emph{todo} $I$ dentro de $I$.

Si en un nodo común el error tiene un patrón $Ch^p$ con términos restantes
menores, la diferencia entre pasos $h$ y $h/2$ permite estimar el error de
la malla fina por $|y_{h/2}-y_h|/(2^p-1)$. Es una estimación condicionada
al patrón, no una cota certificada. Un valor final plausible o dos cifras
impresas iguales no verifican la trayectoria. Recuérdense desigualdades y
potencias. Consulta: solución de \exref{M2-E0311}.

%
\subsection*{Qué se observa y qué se contrasta}
\label{CH03-B03}
Un sensor puede medir un estado, una transformación o una media
$\frac1{b-a}\int_a^b y(t)\,dt$. Escríbase esa relación antes de ajustar
parámetros. Si un dato $d_i$ tiene una cota absoluta garantizada
$\varepsilon_i$, los parámetros admisibles satisfacen
\[
 |H_i(\theta)-d_i|\le\varepsilon_i
 \quad\text{para todos los datos de calibración }i.
\]
Se intersectan estas restricciones con el dominio físico de $\theta$.
Una cota determinista no es una desviación típica.

El dato reservado para contraste se compara con las predicciones de
\emph{todo} el conjunto compatible. Usarlo para recalibrar elimina su
independencia. Sepárense el error numérico, la incertidumbre de medida y
la discrepancia entre modelo y realidad: un rechazo se refiere al conjunto
de hipótesis utilizado. Recuérdense intervalos, logaritmos y media integral.
Consulta: solución de \exref{M2-E0312}.

%
\subsection*{Comparar entradas mediante integrales}
\label{CH03-B04}
La fórmula lineal del capítulo 2 sigue siendo válida, por tramos, para
entradas continuas a trozos y estado continuo. Con datos iniciales iguales,
la diferencia de respuestas a $q_1,q_2$ es
\[
 y_1(t)-y_2(t)=\int_{t_0}^{t}
 \exp\!\left(-\int_s^t p(r)\,dr\right)
 [q_1(s)-q_2(s)]\,ds.
\]
Para $t\ge t_0$, el peso es positivo. El signo del integrando permite
comparar respuestas cuando las entradas están ordenadas; otras
restricciones sobre ellas requieren un argumento adicional.

Una entrada admisible da un ejemplo, no una cota universal. Para establecer
un extremo, pruébese primero la cota y después que se alcanza. Si el
conjunto de entradas es convexo y el observable depende linealmente de
ellas, las combinaciones convexas permiten obtener valores intermedios.
Compruébese cada hipótesis. Consulta: solución de \exref{M2-E0375}.
